\origpage{184}{ROBERTSON AND SEYMOUR}
\noindent is the $(\mathscr{T}_i, \mathscr{T}_j)$-distinction. Let $P$ be the minimal $H_i - H_j$ path in $T$. Then $e_1 \in E(P)$, and since $t_0 \notin H_i \cup H_j$, $E(P)$ contains one of $e_2, \dots, e_k$, say $e_2$. Now
\begin{equation*}
\left( \bigcup_{t \in V(T_2)} \tau(t), \bigcup_{t \in V(T'_2)} \tau(t) \right)
\end{equation*}
has $\lambda$-order strictly less than that of the $(\mathscr{T}_i, \mathscr{T}_j)$-distinction and hence has order at most that of the $(\mathscr{T}_i, \mathscr{T}_j)$-distinction. By (5), $e_2$ makes a separation which distinguishes $\mathscr{T}_i$ from $\mathscr{T}_j$, with $\lambda$-order strictly smaller than that of the $(\mathscr{T}_i, \mathscr{T}_j)$-distinction, a contradiction.

Let $H_i = \{t_i\}$ ($1 \le i \le n$); then the theorem is satisfied. $\square$

We call the tree-decomposition of (10.3) the \textit{standard tree-decomposition} of $G$ relative to $\mathscr{T}_1, \dots, \mathscr{T}_n$.

From (10.3) we deduce a corollary mentioned earlier. We merely sketch the proof since we do not need the result.

\stmt
(10.4) \textit{In any hypergraph $G$ there are at most $|V(G)|$ maximal tangles.}

\textit{Proof.} Let $\mathscr{T}_1, \dots, \mathscr{T}_n$ be the distinct maximal tangles in $G$, and let $\lambda$ be a tie-breaker. Since they are mutually distinguishable, there is a standard tree-decomposition $(T, \tau)$. Let $e, f \in E(T)$ be distinct, making separations $(A, B)$ and $(C, D)$, say, where $A \subseteq C$ and $D \subseteq B$. If $V(A) = V(C)$ then it follows easily that $A = C$, $B = D$, a contradiction; thus $V(A) \subset V(C)$ and similarly $V(D) \subset V(B)$. From this one can show that $|E(T)| \le |V(G)| - 1$, and hence $n = |V(T)| \le |V(G)|$, as required. \fix{We give the details. There are only finitely many tangles in $G$, since each of them is a set of separations of $G$. Two distinct maximal tangles are distinguishable, since if one of them were a subset of the other they would be equal, by maximality. A tie-breaker exists by (9.2), and so if $n \ge 1$, the standard tree-decomposition $(T, \tau)$ exists by (10.3), where $t_i$ corresponds to $\mathscr{T}_i$. Let $\theta_i$ be the order of $\mathscr{T}_i$. For an edge $e$ of $T$ with ends $t_i, t_j$, let $G^e_i$ be the union of $\tau(t)$ over the vertices $t$ of the component of $T \setminus e$ that contains $t_i$, as in the proof of (9.1). By (10.3)(ii), applied to $t_i$ and $t_j$, whose path in $T$ has only the edge $e$, one of the two separations made by $e$ is the $(\mathscr{T}_j, \mathscr{T}_i)$-distinction. It is in $\mathscr{T}_j$ and its reverse is in $\mathscr{T}_i$. It is not $(G^e_j, G^e_i)$, since $(G^e_i, G^e_j) \notin \mathscr{T}_i$ by (10.3)(i). Hence
\begin{equation*}
(G^e_i, G^e_j) \in \mathscr{T}_j \quad \text{and} \quad (G^e_j, G^e_i) \in \mathscr{T}_i.
\end{equation*}

We first prove the inclusions $V(A) \subset V(C)$ and $V(D) \subset V(B)$ of the sketch. As in the proof of (9.1), we may name the ends of $e$ as $t_i, t_j$ and those of $f$ as $t_k, t_l$ so that the ends of $f$ are in the component of $T \setminus e$ that contains $t_j$, and the ends of $e$ are in the component of $T \setminus f$ that contains $t_k$. Then $(A, B) = (G^e_i, G^e_j)$ and $(C, D) = (G^f_k, G^f_l)$ satisfy $A \subseteq C$ and $D \subseteq B$, as shown there. Suppose that $V(A) = V(C)$. Then $V(C \cap D) \subseteq V(A \cap B)$ since $V(D) \subseteq V(B)$, and so $(C, D)$ has order less than $\theta_j$, because $(A, B) \in \mathscr{T}_j$. Since $V(C) \subseteq V(A)$, it follows that $(C, D) \in \mathscr{T}_j$ by (2.9)(ii). But $t_j$ is in the component of $T \setminus f$ that contains $t_k$, and so $(C, D) \notin \mathscr{T}_j$ by (10.3)(i), a contradiction. Thus $V(A) \subset V(C)$. The same argument, applied to $(D, C) \in \mathscr{T}_k$ and $(B, A)$, shows that $V(D) \subset V(B)$.

For the bound on $n$ we show the following.

\hspace*{12pt}(1) \textit{If $t_i, t_j$ are adjacent in $T$, then $V(\tau(t_i)) \not\subseteq V(\tau(t_j))$.}

Let $e$ be the edge joining $t_i$ and $t_j$, and put $(A, B) = (G^e_i, G^e_j)$. Then $(B, A) \in \mathscr{T}_i$, and so $(A, B) \notin \mathscr{T}_i$ by (2.1), and $(A, B)$ has order less than $\theta_i$. Suppose that $V(\tau(t_i)) \subseteq V(\tau(t_j))$. Since $\tau(t_i) \subseteq A$ and $\tau(t_j) \subseteq B$, it follows that $V(\tau(t_i)) \subseteq V(A \cap B)$. Let $t_{k_1}, \dots, t_{k_m}$ be the neighbours of $t_i$ other than $t_j$, let $e_h$ be the edge joining $t_i$ and $t_{k_h}$, and put $(A_h, B_h) = (G^{e_h}_{k_h}, G^{e_h}_i)$ $(1 \le h \le m)$. Then $(A_h, B_h) \in \mathscr{T}_i$, as shown above. Also $V(A_h \cap B_h) \subseteq V(\tau(t_i))$, by (iii) of the definition of tree-decomposition, since $t_i$ is on the path of $T$ between any vertex of the component of $T \setminus e_h$ that contains $t_{k_h}$ and any vertex of the other component. If $m = 0$, then $t_i$ is the only vertex of its component of $T \setminus e$, so $A = \tau(t_i)$ and $V(A) \subseteq V(B)$. Then $V(B) = V(G)$, contrary to the third axiom, since $(B, A) \in \mathscr{T}_i$. So $m \ge 1$. For $1 \le h \le m$, every vertex of $(A_1 \cup \dots \cup A_h) \cap (B_1 \cap \dots \cap B_h)$ is in $A_g \cap B_g$ for some $g \le h$, and so it is in $V(\tau(t_i)) \subseteq V(A \cap B)$. Hence this separation has order less than $\theta_i$, and by (2.2) and induction on $h$ it is in $\mathscr{T}_i$. Put $C = A_1 \cup \dots \cup A_m$ and $D = B_1 \cap \dots \cap B_m$, so that $(C, D) \in \mathscr{T}_i$. We claim that $V(D) \subseteq V(B)$. Let $x \in V(D)$, and suppose that $x \notin V(B)$. Then $x \in V(A)$, and so $x \in V(\tau(t))$ for some vertex $t$ of the component of $T \setminus e$ that contains $t_i$. Here $t \ne t_i$ since $V(\tau(t_i)) \subseteq V(B)$, and so the path of $T$ from $t_i$ to $t$ has first edge $e_h$ for some $h$. Since $x \in V(B_h)$, there is a vertex $t'$ of the component of $T \setminus e_h$ that contains $t_i$ with $x \in V(\tau(t'))$. Then $t_i$ is on the path of $T$ between $t'$ and $t$, and so $x \in V(\tau(t_i)) \subseteq V(B)$, a contradiction. Thus $V(D) \subseteq V(B)$, and so $(A, B) \in \mathscr{T}_i$ by (2.9)(i), a contradiction. This proves (1).

If $n = 0$ there is nothing to prove, and if $n = 1$ then $n \le |V(G)|$ since $\mathscr{T}_1$ has order at least 1 and at most $|V(G)|$. So we may assume that $n \ge 2$. For each vertex $t \ne t_1$ of $T$, let $p(t)$ be the neighbour of $t$ on the path of $T$ from $t$ to $t_1$, let $e_t$ be the edge joining $t$ and $p(t)$, and choose $x(t) \in V(\tau(t)) - V(\tau(p(t)))$, which is possible by (1). Since $n \ge 2$, $t_1$ has a neighbour, and so $V(\tau(t_1)) \ne \emptyset$ by (1). Choose $x(t_1) \in V(\tau(t_1))$. We claim that $x(t) \ne x(t')$ for all distinct $t, t'$. Suppose that $x(t) = x(t') = x$. If $t \ne t_1$ and $t'$ is in the component of $T \setminus e_t$ that contains $p(t)$, then $p(t)$ is on the path of $T$ between $t$ and $t'$, and so $x \in V(\tau(p(t)))$ by (iii), contrary to the choice of $x(t)$. Similarly it is impossible that $t' \ne t_1$ and $t$ is in the component of $T \setminus e_{t'}$ that contains $p(t')$. If $t = t_1$, then $t' \ne t_1$, and $t = t_1$ is in the component of $T \setminus e_{t'}$ that contains $p(t')$, by the choice of $p(t')$, which we have seen is impossible. So $t \ne t_1$, and similarly $t' \ne t_1$. Hence $t'$ is in the component of $T \setminus e_t$ that contains $t$, and $t$ is in the component of $T \setminus e_{t'}$ that contains $t'$. In particular $e_t \ne e_{t'}$, since otherwise $t' = p(t)$. The component of $T \setminus e_t$ that contains $p(t)$ also contains $t_1$, by the choice of $p(t)$, and it does not contain $t'$. So it does not contain $e_{t'}$, and $p(t)$ and $t_1$ are in the same component of $T \setminus e_{t'}$. Since $t'$ and $t$ are in the same component of $T \setminus e_{t'}$, and $t$ and $p(t)$ are joined by $e_t \ne e_{t'}$, it follows that $t'$ and $t_1$ are in the same component of $T \setminus e_{t'}$. This is contrary to the choice of $p(t')$, which puts $t'$ and $t_1$ in different components of $T \setminus e_{t'}$. This proves our claim. Hence $n = |V(T)| \le |V(G)|$, and $|E(T)| = n - 1 \le |V(G)| - 1$.}\fixnote{fix:10.4-proof-added}{The paper only sketches this proof. The details above are the argument used in the formalization.} $\square$

\sectionhead{11. \textsc{Structure} \textsc{Relative} \textsc{To} A \textsc{Tangle}}

Now we come to the last main result of the paper. We have seen in (5.2) that if $G$ has small tangle number, then it has a tree-decomposition of small width. Our problem here is, suppose that $G$ has large tangle number, but relative to each high order tangle the graph has a structure or decomposition of a certain kind $X$, say; what can we infer about the global structure of $G$ from this local knowledge? One might guess that $G$ should have a tree-decomposition into pieces each with structure $X$, but that is false. Nevertheless, it turns out that $G$ has a tree-decomposition into pieces which ``almost'' have structure $X$, and we need to know this for an application in [6].
